\section{Introduction}\label{sec:intro}
Autonomous vehicles rely on perception models, which are typically neural networks that interpret sensor measurements in order to detect, classify and localise objects in the surrounding scene. For LiDAR, a common perception task is semantic segmentation, in which every point of a scan is assigned a class such as car, road or pedestrian. Training such models requires large and varied datasets in which every point is labelled, and recording and annotating these datasets is a costly and time-consuming process~\cite{xiaoTransferLearningSynthetic2022}. The annotation of the SemanticKITTI benchmark alone, for example, required more than 1700 hours of work~\cite{behleySemanticKITTIDatasetSemantic2019a}. Synthetic data generated in simulation offers an alternative, since the simulator provides control over the content of each scene and can label every point automatically~\cite{hurlPreciseSyntheticImage2019,xiaoTransferLearningSynthetic2022}. Synthetic data nevertheless differs from real sensor measurements, and this difference, known as the domain gap, can limit how well a model trained on synthetic data performs on real data~\cite{hurlPreciseSyntheticImage2019,manivasagamLiDARsimRealisticLiDAR2020a}.

Previous work shows that synthetic data can still be valuable despite this gap. Our literature review found that synthetic datasets and simulators such as PreSIL, LiDARsim and SynLiDAR yield useful transfer to real data, although only in particular settings~\cite{hurlPreciseSyntheticImage2019,manivasagamLiDARsimRealisticLiDAR2020a,xiaoTransferLearningSynthetic2022}. The benefit was mainly reported when the synthetic scans were used for pre-training followed by fine-tuning on real scans~\cite{hurlPreciseSyntheticImage2019}, when the simulator reproduced real sensor effects such as missing returns~\cite{manivasagamLiDARsimRealisticLiDAR2020a}, or when domain adaptation was applied to reduce the gap~\cite{xiaoTransferLearningSynthetic2022}. Because the benefit depends on the setting, the value of a synthetic dataset cannot be assumed in advance and has to be established for each case. The review further identified the performance on held-out real data of a model trained on the synthetic data, referred to in this proposal as downstream performance, as the most direct evidence of this value.

Obtaining this evidence is costly, as a model has to be trained and evaluated for every candidate dataset or sensor configuration. For this reason, generated scans are commonly assessed with data-level diagnostics, which compare a set of synthetic scans directly with a set of real scans without training a model~\cite{zyrianovLearningGenerateRealistic2022}. This proposal considers three such diagnostics, which were used to evaluate the generative model LiDARGen~\cite{zyrianovLearningGenerateRealistic2022}. Each diagnostic summarises the comparison in a single value, and a lower value indicates that the synthetic and real sets are more similar. The \emph{Fr\'echet range distance} (FRD) passes both sets through a pretrained segmentation network, RangeNet++~\cite{miliotoRangeNet++Fast2019a}, and compares the mean and covariance of the extracted features, thereby measuring whether the network responds to both sets in a similar way. The \emph{maximum mean discrepancy} (MMD) and the \emph{Jensen--Shannon divergence} (JSD) instead compare where the points are located around the sensor. Both are computed on bird's-eye-view (BEV) maps, which divide the area around the sensor into a top-down grid and count the number of points in each cell. MMD compares the individual maps of the real set with those of the synthetic set, whereas JSD first sums the maps of each set and then compares the two resulting distributions of points over the grid.

In this proposal, downstream performance is measured with the mean intersection over union (mIoU). For each class, the intersection over union is the overlap between the predicted and true labels divided by their union, and the mIoU is the average of this value over all classes. In contrast to the three diagnostics, which summarise an entire set of scans in a single value, the mIoU can also be reported per class and per distance range. This makes it possible to detect a loss of performance on specific cases, such as small or distant objects, that a value computed over complete scans may fail to reflect.

Whether the diagnostics can replace real-data evaluation thus depends on how closely they follow changes in downstream performance. To our knowledge, no study has quantified how far synthetic data must deviate from real data before the diagnostics register the deviation, nor how this compares with the deviation at which downstream performance starts to decline. The proposed study addresses this gap with a controlled experiment in which ray drop is applied to a synthetic dataset in increasing steps and the response of both the diagnostics and the mIoU is measured at every step. Ray drop refers to emitted laser pulses for which no return is measured, an effect that occurs in real sensors but is absent from basic simulations~\cite{manivasagamLiDARsimRealisticLiDAR2020a}. The results are expected to show to what extent a favourable diagnostic score justifies omitting real-data evaluation. Such insight can reduce the cost of developing synthetic LiDAR datasets without overlooking failures that matter for perception.